% Capítulo 3 (Principais Trabalhos Relacionados) — Slide 1 — Panorama da literatura
\begin{frame}{Panorama da literatura}
  \centering
  {\footnotesize
  \renewcommand{\arraystretch}{1.15}
  \begin{tabular}{@{}p{3.0cm}p{3.6cm}p{4.6cm}@{}}
    \toprule
    \textbf{Família} & \textbf{Exemplos} & \textbf{Limitação} \\
    \midrule
    Políticas clássicas & EOQ, $(s,S)$, Jornaleiro & Baixa adaptação a séries heterogêneas \\
    Previsão de demanda & Croston, ARIMA, LSTM, XGBoost & Não define a política a aplicar \\
    Parametrização por meta-heurísticas & GA, SA, PSO, DE, MILP & Calibra política já escolhida, não seleciona entre famílias \\
    Aprend.\ por reforço & DQN, PPO, SARSA & Exige históricos longos ($\sim 10^6$ passos) \\
    Seleção/meta-aprend. & ASP, hiper-heurísticas, AutoML & Pouco explorado em varejo real \\
    \bottomrule
  \end{tabular}}
  \vspace{0.6em}

  \small \fonte{capítulo de Trabalhos Relacionados, Tabela 3.1}
  \note{
    A literatura oferece blocos sólidos para cada parte do problema, mas
    isolados. As políticas clássicas são simples e interpretáveis, mas se
    adaptam mal a séries heterogêneas. A previsão de demanda intermitente
    melhora estimativas, mas não decide qual política usar. A parametrização
    por meta-heurísticas calibra bem uma política já escolhida, mas não
    responde qual família adotar. Agentes de aprendizado por reforço aprendem
    decisões sequenciais,
    mas exigem históricos da ordem de um milhão de interações -- incompatível
    com séries de varejo de 38 ciclos. E a literatura de seleção de
    algoritmos e meta-aprendizado, embora madura em outras áreas, é pouco
    explorada para selecionar políticas de reposição em redes varejistas
    reais. É exatamente nesse espaço entre os blocos que o AIPE se posiciona.
  }
\end{frame}
